
\section{The specification surface}
\label{sec:surface}

\subsection{The design space}

\begin{definition}[Design space]
A \emph{design space} $\mathcal{S}$ for a feature-value claim is a finite set
of specifications, each a complete assignment to the axes
\[
\begin{aligned}
s \;=\; \bigl(\;
&\underbrace{A}_{\text{learner}},\;
\underbrace{E}_{\text{encoding}},\;
\underbrace{Y}_{\text{target}},\;
\underbrace{\Pi}_{\text{split}},\;
\underbrace{\tau}_{\text{decision time}},\; \\[2pt]
&\underbrace{q}_{\text{register quality}},\;
\underbrace{B}_{\text{baseline}},\;
\underbrace{m}_{\text{metric}},\;
\underbrace{\theta}_{\text{operating point}}\;\bigr),
\end{aligned}
\]
together with the cohort and eligibility rule that fix $D$. A claim about a
feature is \emph{well posed} only relative to a declared $\mathcal{S}$.
\end{definition}

Two of these axes deserve comment because they are usually not treated as
design choices at all.

The \textbf{decision time} $\tau$ fixes which fields exist when the prediction
is made. It is not the same as the split rule and it is not a data-cleaning
detail: moving $\tau$ later in a process admits fields written in between, and
Section~\ref{sec:cmdb} is a case in which that single move takes a field's
increment from resolvably positive to indistinguishable from zero. Feature
availability is an argument of the estimand.

The \textbf{register quality} $q$ acts on $f$ rather than on the comparison. A
field that is a lookup into a maintained register --- a configuration item, a
customer master record, a product catalogue --- inherits that register's
completeness and accuracy. We implement \nQualityKinds\ mechanisms besides the clean
field, each parameterised by the training half alone: population loss with the
long tail absent, with the core absent, and at random; identity error, in
which a share of rows carry a wrong value drawn from the training alphabet;
reconciliation failure, in which a share of identities exist twice; and
staleness, in which an identity is present only if the register had already
seen it by the split point. The first three model coverage, the next two model
accuracy, the last models discovery; the taxonomy follows the event-log
imperfection patterns of \citet{suriadi2017imperfection} where they apply. Each is verified by execution rather than
by assertion to be a function of the training half alone: perturbing the test
half and re-degrading must leave the training half's degraded values
identical, and \texttt{results/s01\_trainonly.csv} records the \nTrainOnly\
checks that enforce it.

\subsection{The surface, and why it is not a scalar}

Given $\mathcal{S}$ and data $D$, the surface is the map $s \mapsto V_s(f)$
of~\eqref{eq:estimand}. A conventional report is one cell of it. The paper's
claim is that a single cell is not a summary of the surface in any useful
sense: Section~\ref{sec:multilog} shows that the sign of $V_s$ disagrees
across cells on most log--target pairs we study, and
Section~\ref{sec:results-regret} measures how often.

There is one state in which a scalar is defensible, and we name it rather than
denying that it exists: when the whole admissible surface lies on the same
side of zero with simultaneous confidence, a positive number is a lower bound
on a robust conclusion. Section~\ref{sec:regions} makes that a definition.

\subsection{The relative reduction, and its pathologies}

Comparisons between two baselines are often stated as a proportion. Write
$B_0 \subset B_1$ and
\begin{equation}
D_s(f) \;=\; V_s(f \mid B_0) - V_s(f \mid B_1),
\quad
R_s(f) \;=\; 1 - \frac{V_s(f \mid B_1)}{V_s(f \mid B_0)}
\;=\; \frac{D_s(f)}{V_s(f \mid B_0)} .
\label{eq:reduction}
\end{equation}
$D$ is an absolute difference in the metric's own units. $R$ is a ratio and
inherits every pathology of one: the denominator can be small, zero or
negative; $R$ is not comparable across metrics, because
Proposition~\ref{prop:nonid} says the two are not functions of each other; and
$R$ can exceed $1$ or fall below $0$, so reading it as ``the share absorbed''
is a category error whenever it leaves $[0,1]$.

We therefore make $D$ and the two increments primary throughout and treat $R$
as a secondary descriptive statistic subject to three rules, declared before
any $R$ in this paper was computed:
\begin{enumerate}
\item $R$ is reported only where $V_s(f \mid B_0)$ is resolvably positive
under the \emph{simultaneous} band, not merely positive as a point estimate;
\item its interval is a Fieller set and its \emph{kind} --- bounded,
unbounded, or the complement of an interval --- is printed beside it, because
a percentile interval reports a bounded interval in all three cases
(Section~\ref{sec:fieller});
\item the magnitude of $R$ is never compared with the spread of $R$ across
metrics as though both lived on an unproblematic common scale.
\end{enumerate}

\subsection{Two propositions and one computational result}
\label{sec:props}

Throughout, a \emph{configuration} is a quadruple of score vectors on a fixed
labelled test set,
$\mathcal{C} = (\sigma_{B_0}, \sigma_{B_0 \cup f}, \sigma_{B_1},
\sigma_{B_1 \cup f})$ with $B_0 \subset B_1$, and $V_m(B)$,
$R_m(\mathcal{C})$ are as in~\eqref{eq:estimand} and~\eqref{eq:reduction}.
Conventions matter here and are fixed: ROC AUC is the Mann--Whitney statistic
with the mid-rank tie rule; average precision is the step-wise sum
$\sum_k (\mathrm{Rec}_k - \mathrm{Rec}_{k-1})\,\mathrm{Prec}_k$ with equal
scores collapsed into one operating point and no interpolation; net benefit is
the Vickers--Elkin per-case quantity with selection rule $p \geq \theta$. The
constructions below depend on these conventions and would need re-tuning under
others; \texttt{results/s07\_conventions.csv} carries the list.

\begin{proposition}[Metric non-identifiability of the reduction]
\label{prop:nonid}
There is no function $\varphi : \mathbb{R} \to \mathbb{R}$ with
$R_{\mathrm{AP}}(\mathcal{C}) =
\varphi\bigl(R_{\mathrm{AUC}}(\mathcal{C})\bigr)$ for every configuration
$\mathcal{C}$ whose denominators are positive under both metrics.
Consequently the reduction is not identified by the data, the feature and the
two baselines: the metric is an argument of the estimand.
\end{proposition}

\begin{proof}[Proof (construction)]
Take $P$ positives and $N$ negatives and let every score take one of two
values, so a scoring rule is described by the pair $(a,b)$ counting positives
and negatives in the high block; write $\alpha = a/P$, $\beta = b/N$,
$\rho = N/P$. Both baselines are the fully tied rule $\alpha = \beta = 1$,
which has $\mathrm{AUC} = 1/2$ and $\mathrm{AP} = \pi$, the prevalence. Then
exactly
\[
\mathrm{AUC} - \tfrac12 = \tfrac{\alpha - \beta}{2},
\qquad
\mathrm{AP} - \pi = \alpha\left(\frac{\alpha}{\alpha + \beta\rho} - \pi\right).
\]
Fixing $\alpha - \beta$ equal on the two legs makes the AUC increments equal,
hence $R_{\mathrm{AUC}} = 0$ exactly --- in integers, not to rounding ---
while sliding $\alpha$ along the admissible range moves the AP increments and
sweeps $R_{\mathrm{AP}}$. With $P = \nPropP$ and $N = \nPropN$ the two
configurations in \texttt{results/s07\_prop1.csv} have identical
$R_{\mathrm{AUC}}$ to \propOneAucGap\ and $R_{\mathrm{AP}}$ differing by
\propOneApGap. A single $\varphi$ cannot take one input to two outputs.
\end{proof}

\begin{remark}
An earlier version of this work called this ``metric-unboundedness''. That is
the wrong word: $R$ is constructed over a bounded interval in any fixed
configuration. What the construction establishes is non-identifiability, and a
\emph{range}: over the sweep in \texttt{s07\_prop1.csv}, $R_{\mathrm{AUC}}$
stays at zero to \propOneSweepAuc\ while $R_{\mathrm{AP}}$ reaches
\propOneSweepAp.
\end{remark}

\begin{proposition}[Rank invariance]
\label{prop:rank}
Let $\psi$ be strictly increasing and continuous, applied
\emph{simultaneously} to all four score vectors of a configuration --- one
recalibration of the scoring system, not a different one per arm --- and call
$m$ \emph{rank-based} if $m(\psi(\sigma), y) = m(\sigma, y)$ for every such
$\psi$ and every $(\sigma, y)$.
\begin{enumerate}
\item[(i)] If $m$ is rank-based then $V_m$ and $R_m$ are invariant under
$\psi$, for every configuration and every $\psi$, wherever $R_m$'s denominator
is nonzero.
\item[(ii)] If $m$ is not rank-based then there \emph{exist} a configuration
and a $\psi$ with $R_m(\psi\mathcal{C}) \neq R_m(\mathcal{C})$.
\end{enumerate}
\end{proposition}

\begin{proof}
(i) is immediate: $V_m$ is a difference of two invariants and $R_m$ a ratio of
two such differences, so both are unchanged wherever they are defined. (ii) is
existential and is witnessed by construction. On the configuration of
\texttt{results/s07\_prop2.csv}, with $\psi(p) = p^2$, the rank-based
instruments move $R$ by \propTwoRankShift\ and the non-rank-based ones by up
to \propTwoNonRankShift.
\end{proof}

\begin{remark}[What the necessity direction cannot say]
\label{rem:cancel}
The quantifier in (ii) is existential, and the paper now says so. The
universal reading --- that a non-rank-based metric \emph{always} moves $R$ ---
is false, because the change in the two legs can cancel in the ratio. The same
file exhibits such a configuration: a non-rank-based instrument whose $V$
moves under $\psi$ and whose $R$ does not, to \propTwoCancelShift. An earlier
version of this work asserted the universal reading; that is
correction~C13 in \ref{app:corrections}.
\end{remark}

\begin{result}[Axis non-redundancy on a finite family]
\label{res:nonredundancy}
For an ordered pair of axes $(i,j)$, call a \emph{witness} a pair of design
points that agree on every axis except $i$, give the same increment at the
reference level of axis $j$, and give different increments at another level of
$j$. A witness shows that knowing the answer along axis $i$ does not determine
the answer along axis $j$. Over a family of \nResultEstates\ synthetic estates
and the \nResultPairs\ ordered pairs of the four reporting axes, a witness is
found for between \resultWitnessMin\ and \resultWitnessMax\ pairs as the
equality tolerance is swept (\texttt{results/s07\_result1.csv}).
\end{result}

\begin{remark}
This is not a proposition and is not labelled as one. It is a finite search;
its answer moves with an arbitrary tolerance, which is why the whole sweep is
printed; failing to find a witness is not evidence that one axis determines
another; and finding witnesses in a family of \nResultEstates\ estates is not
a general non-redundancy theorem. An earlier version of this work presented it
as a proposition, which was the third of the referee's objections to the
theory section, and it is granted.
\end{remark}
